\documentclass[conference]{IEEEtran}

\usepackage{graphicx}
\usepackage{amsmath,amssymb}
\usepackage{float}
\usepackage{booktabs}
\usepackage{multirow}

\begin{document}

\title{
    \textbf{Benchmark Report: 557.xz\_r} \\
    \text{Register Spill Analysis on RISC-V via gem5}
}

\author{
    \IEEEauthorblockN{Lütfullah Burak Kaya}
    \IEEEauthorblockA{
        Ozyegin University \\
        burak.kaya.50711@ozu.edu.tr
    }
    \and
    \IEEEauthorblockN{Ismail Akturk}
    \IEEEauthorblockA{
        Ozyegin University \\
        ismail.akturk@ozyegin.edu.tr
    }
}

\newcommand{\LongDate}{%
    \number\day\space
    \ifcase\month
    \or January\or February\or March\or April\or May\or June%
    \or July\or August\or September\or October\or November\or December%
    \fi
    \space \number\year
}

\twocolumn[
\begin{@twocolumnfalse}
\maketitle
\begin{center}{\small \LongDate}\end{center}
\vspace{1em}
\end{@twocolumnfalse}
]

% ================================================

\begin{abstract}
This report presents the register spill characterization of the
\texttt{557.xz\_r} benchmark from SPEC CPU2017, executed on a
RISC-V (RV64GC) target using the gem5 TimingSimpleCPU simulator.
A custom SpillDetector was integrated into gem5 to count store-then-load
patterns on the stack within a user-defined Region of Interest (ROI).
The test dataset (compression of \texttt{cpu2006docs.tar.xz} at
compression level~0) was used as the input. Results show a spill rate
of \textbf{1.91\%} within the ROI, corresponding to 38.4 million
detected spill events across 2.01 billion ROI instructions. The LZMA
compression algorithm's predominantly sequential data-processing
pipeline produces relatively low register pressure compared to the
recursive or tensor-field workloads in the evaluated suite.
\end{abstract}

% =========================================================
\section{Benchmark Overview}

\texttt{557.xz\_r} is the SPEC CPU2017 rate variant of the XZ
compression utility (liblzma), implementing the LZMA2 compression
algorithm. The benchmark compresses and decompresses data archives
in memory, measuring throughput across a range of compression levels
and dictionary sizes. The LZMA encoder maintains a sliding window
dictionary, a match-finding data structure (hash chains or binary
trees), and a range encoder with probability tables; the decoder
reverses this process via a range decoder and a copy-on-match step.

The RISC-V binary (\texttt{xz\_r.riscv}) was compiled with gem5
ROI markers inserted in \texttt{src/spec.c}, wrapping the main
compression for-loop:

\begin{itemize}
    \item \texttt{m5\_work\_begin(0,0)} --- before the compression loop
    \item \texttt{m5\_work\_end(0,0)}   --- after the loop completes
\end{itemize}

The ROI encompasses only the compression phase; decompression of the
input archive (used to reconstruct the source data in memory) occurs
before the ROI begins.

\subsection{Input Datasets}

SPEC CPU2017 provides three dataset sizes for \texttt{xz\_r},
differing in input file size, number of threads, and the set of
compression levels tested (Table~\ref{tab:inputs}).

\begin{table}[H]
\centering
\caption{xz\_r Input Dataset Sizes}
\label{tab:inputs}
\begin{tabular}{lll}
\toprule
\textbf{Dataset} & \textbf{Input archive} & \textbf{Levels tested} \\
\midrule
test    & cpu2006docs.tar (8.6\,MiB)  & 6 levels (0--4e) \\
train   & Similar archives            & Extended level set \\
refrate & cld.tar + others (multi-GB) & Many levels, repeated \\
\bottomrule
\end{tabular}
\end{table}

The \textbf{test} dataset was selected for this study. The SPEC harness
decompresses \texttt{cpu2006docs.tar.xz} (1.3\,MiB on disk, 8.6\,MiB
uncompressed) into a memory buffer before the ROI, then the binary
re-compresses the data with the specified parameters. For this
simulation, the first control file entry was used: compression with
4 threads at level~0 (fastest compression), producing an output
between 1{,}548{,}636 and 1{,}555{,}348 bytes as verified by the
benchmark harness.

% =========================================================
\section{Simulation Setup}

\subsection{gem5 Configuration}

\begin{itemize}
    \item \textbf{CPU model:}   TimingSimpleCPU
    \item \textbf{ISA:}         RISC-V RV64GC
    \item \textbf{Caches:}      enabled (default L1 I/D + L2)
    \item \textbf{Memory:}      4\,GB simulated physical RAM
    \item \textbf{Spill mode:}  summary only (\texttt{spill\_verbose=False})
\end{itemize}

The full gem5 invocation was:

\begin{verbatim}
./build/RISCV/gem5.opt
  --outdir=benchmarks/xz/m5out_test
  configs/deprecated/example/se.py
  --cmd=benchmarks/xz/xz_r.riscv
  --options="benchmarks/xz/cpu2006docs.tar.xz
             4 <sha512> 1548636 1555348 0"
  --cpu-type=TimingSimpleCPU
  --caches
  --mem-size=4GB
\end{verbatim}

% =========================================================
\section{Simulation Results}

The simulation completed successfully. The benchmark harness verified
that the compressed output size fell within the expected range,
confirming a correct compression result.

\subsection{Pre-ROI Context}

Before the compression ROI, the binary decompressed the input
\texttt{cpu2006docs.tar.xz} archive into a memory buffer. This
decompression phase is substantial, accounting for 932 million
instructions and representing 31.7\% of total simulated instructions.
Table~\ref{tab:preroi} shows the global counters at ROI entry.

\begin{table}[H]
\centering
\caption{Global Counters at ROI Entry}
\label{tab:preroi}
\begin{tabular}{lr}
\toprule
\textbf{Metric} & \textbf{Value} \\
\midrule
Total instructions (pre-ROI) & 932,903,724 \\
Total spills detected (pre-ROI) & 16,591,613 \\
Pre-ROI spill rate & 1.78\% \\
\bottomrule
\end{tabular}
\end{table}

The pre-ROI spill rate of 1.78\% is consistent with the ROI spill
rate (1.91\%), reflecting that both the decompression and compression
phases of XZ share similar sequential data-processing characteristics
and produce comparable register pressure.

\subsection{ROI Spill Statistics}

Table~\ref{tab:roi} presents the spill statistics collected within
the ROI.

\begin{table}[H]
\centering
\caption{ROI Spill Statistics --- 557.xz\_r (test, level 0)}
\label{tab:roi}
\begin{tabular}{lr}
\toprule
\textbf{Metric} & \textbf{Value} \\
\midrule
ROI instructions         & 2,014,351,522 \\
ROI stores               &   183,548,298 \\
ROI loads                &   367,733,280 \\
\midrule
\textbf{ROI spills}      & \textbf{38,411,428} \\
\textbf{Spill rate}      & \textbf{1.9069\%} \\
\bottomrule
\end{tabular}
\end{table}

\subsection{Timing}

\begin{table}[H]
\centering
\caption{Simulation Timing}
\label{tab:timing}
\begin{tabular}{lr}
\toprule
\textbf{Metric} & \textbf{Value} \\
\midrule
Simulated time           & 5.17\,s \\
Host wall time           & 1{,}606.16\,s ($\approx$26.8\,min) \\
Total simulated instructions & 2,947,406,976 \\
\bottomrule
\end{tabular}
\end{table}

% =========================================================
\section{Analysis}

\subsection{Spill Rate Interpretation}

The measured spill rate of \textbf{1.91\%} places \texttt{xz\_r}
twelfth among all evaluated SPEC CPU2017 benchmarks
(Table~\ref{tab:compare}), in the low-pressure tier alongside
\texttt{538.imagick\_r} (2.06\%) and \texttt{505.mcf\_r} (1.83\%).

\begin{table}[H]
\centering
\caption{Cross-Benchmark Spill Rate Comparison}
\label{tab:compare}
\begin{tabular}{lrr}
\toprule
\textbf{Benchmark} & \textbf{ROI Spills} & \textbf{Spill Rate} \\
\midrule
507.cactuBSSN\_r & 3,782,099,438 & 7.8756\% \\
511.povray\_r    &   174,359,922 & 7.7809\% \\
500.perlbench\_r &     1,070,055 & 7.2838\% \\
502.gcc\_r       &     1,108,838 & 6.9185\% \\
526.blender\_r   &    61,845,869 & 6.2193\% \\
523.xalancbmk\_r &    18,381,250 & 4.6427\% \\
541.leela\_r     & 1,158,236,271 & 4.4438\% \\
519.lbm\_r       &   274,130,440 & 4.2130\% \\
531.deepsjeng\_r & 1,928,903,656 & 4.1749\% \\
544.nab\_r       &   151,310,616 & 2.2612\% \\
538.imagick\_r   &     2,424,899 & 2.0601\% \\
557.xz\_r        &    38,411,428 & \textbf{1.9069\%} \\
505.mcf\_r       &   446,113,136 & 1.8263\% \\
508.namd\_r      &   327,718,955 & 1.0269\% \\
\bottomrule
\end{tabular}
\end{table}

\subsection{Store vs.\ Load Balance}

The load count (367.7M) is $2.00\times$ higher than the store count
(183.5M), a near-perfect 2:1 ratio. This symmetric balance reflects
LZMA's sliding-window architecture: for every byte written to the
compressed output stream, the encoder reads the corresponding source
byte and one or more dictionary entries to find the best match.
The match-finder updates its hash table (one write) and probes it
(multiple reads), while the range encoder alternately reads and writes
probability tables in a regular pattern. The clean 2:1 ratio indicates
a structured, predictable memory access pattern unlike the irregular
traversals seen in tree-search or interpreter workloads.

\subsection{Spill Fraction of Loads}

Within the ROI, the ratio of spills to total loads is:
\[
\frac{38{,}411{,}428}{367{,}733{,}280} \approx 10.4\%
\]
Approximately 1 in 10 load instructions is a spill reload. This is
one of the lowest spill-to-load fractions in the evaluated suite.
LZMA compression's inner loop maintains a compact working set: the
current position in the input stream, the match-finder state, and
the range encoder state. The RISC-V backend can keep most of these
in registers for extended stretches of the inner loop, spilling only
when entering the less-frequently-called outer functions.

\subsection{Pre-ROI Spill Density}

The pre-ROI phase (decompression, 932M instructions, 1.78\% spill
rate) is notably large, accounting for 31.7\% of total simulated
instructions. This makes \texttt{xz\_r} the benchmark with the
largest absolute pre-ROI instruction count in the study. The
decompression phase shows nearly identical register pressure to
the compression ROI (1.78\% vs.\ 1.91\%), confirming that both
LZMA encoding and decoding share structurally similar sequential
data-flow patterns.

\subsection{ROI Coverage}

The ROI accounts for $2{,}014{,}351{,}522$ out of $2{,}947{,}406{,}976$
total simulated instructions:
\[
\frac{2{,}014{,}351{,}522}{2{,}947{,}406{,}976} \approx 68.3\%
\]
The 68.3\% ROI coverage is the lowest of all evaluated benchmarks,
a direct consequence of the large decompression pre-ROI. However,
this is by design: the benchmark must reconstruct the input data
in memory before compressing it. The ROI correctly isolates the
compression phase, which is the intended measured workload.

\subsection{LZMA Sequential Data Flow}

The low spill rate of 1.91\% reflects the structured, sequential
nature of LZMA compression. Unlike recursive workloads
(\texttt{povray\_r}, \texttt{deepsjeng\_r}) or tensor-field
computations (\texttt{cactuBSSN\_r}), LZMA's inner loop operates on
a fixed-size state that fits comfortably within RISC-V's 32 registers:
the input pointer, the match-finder context, the range encoder state,
and a small number of probability table pointers. The near-constant
1.78--1.91\% spill rate across both decompression and compression
phases confirms that register pressure in \texttt{xz\_r} is
determined by the algorithm's inherent data-flow structure rather
than by the specific compression level or input content.

% =========================================================
\section{Conclusion}

The \texttt{557.xz\_r} benchmark exhibits a register spill rate of
1.91\% within the compression ROI on RISC-V RV64GC, placing it
twelfth among the 14 evaluated SPEC CPU2017 benchmarks in the
low-pressure tier. With 38.4 million spill events across 2.01 billion
ROI instructions and a near-symmetric 2.00:1 load-to-store ratio,
\texttt{xz\_r} presents the most structured and predictable memory
access pattern of the evaluated suite. The large pre-ROI
decompression phase (932M instructions, 1.78\% spill rate) mirrors
the ROI characteristics, confirming the structural regularity of
LZMA data flow. These results demonstrate that sequential
data-compression workloads are significantly less register-pressure
intensive on RISC-V RV64GC than recursive, object-oriented, or
physics-simulation workloads.

\end{document}
